\documentclass[11pt,a4paper]{article}
\usepackage[T1]{fontenc}
\usepackage[utf8]{inputenc}
\usepackage[brazil]{babel}
\usepackage{amsmath,amssymb}
\usepackage{geometry}
\usepackage{hyperref}
\usepackage{booktabs}
\usepackage{enumitem}
\usepackage{xcolor}
\geometry{margin=1in}

\hypersetup{
  colorlinks=true,
  linkcolor=blue,
  urlcolor=blue,
  citecolor=blue,
  pdftitle={Softmax, Ajuste de Pesos e Classe Sem Transtorno},
  pdfauthor={Assistente},
  pdfcreator={Assistente}
}

\title{\textbf{Softmax, Ajuste de Pesos e Classe \emph{Sem Transtorno}}}
\author{Guia didático e objetivo}
\date{\today}

\begin{document}
\maketitle

\section{Softmax: do escore à probabilidade relativa}
Considere um vetor de escores (também chamados de \emph{logits}) $s=(s_1,\ldots,s_K) \in \mathbb{R}^K$.
A função \textbf{softmax} produz uma distribuição de probabilidades $p=(p_1,\ldots,p_K)$ tal que $p_k \in (0,1)$ e $\sum_k p_k = 1$:
\begin{equation}
  \mathrm{softmax}(s)_k \;=\; \frac{e^{s_k}}{\sum_{t=1}^{K} e^{s_t}} \,.
\end{equation}
Para estabilidade numérica, subtrai-se $\max(s)$ (não altera o valor final):
\begin{equation}
  \mathrm{softmax}(s)_k \;=\; \frac{e^{\,s_k - \max(s)}}{\sum_{t=1}^{K} e^{\,s_t - \max(s)}} \,.
\end{equation}

\paragraph{Intuição.} As exponenciais ampliam diferenças: uma pequena vantagem em $s_k$ pode virar uma vantagem maior em $p_k$.
A normalização pela soma garante que os $p_k$ somem $1$, permitindo comparar classes em uma mesma linha.

\paragraph{Temperatura (opcional).} Controla a ``nitidez'' das probabilidades:
\begin{equation}
  \mathrm{softmax}_T(s)_k \;=\; \frac{e^{s_k/T}}{\sum_t e^{s_t/T}}\,,
\end{equation}
com $T<1$ deixando a distribuição mais ``dura'' (top-1 maior) e $T>1$ mais ``suave''.

\section{Ajuste de pesos: como melhorar sem ``estragar'' o questionário}
Seja $X \in [0,1]^{n \times m}$ a matriz de respostas (linhas = casos, colunas = itens/perguntas) e $W \in [0,1]^{m \times K}$ a matriz de pesos (colunas = classes).
Calculamos notas lineares $S$ e probabilidades $P$ por:
\begin{equation}
  S \;=\; XW \quad\text{e}\quad P_{i,:} \;=\; \mathrm{softmax}(S_{i,:}) \,.
\end{equation}

\subsection*{Objetivo e métrica}
Nosso objetivo é aumentar a taxa de acerto por classe no \textbf{top-3} (``macro top-3''), isto é, a média (entre classes) da fração de linhas em que a classe verdadeira aparece entre as 3 maiores probabilidades.
Para suportar alvos com \emph{múltiplos rótulos}, definimos uma distribuição-alvo por linha $Y \in [0,1]^{n \times K}$ onde, para os rótulos presentes na linha $i$, atribuimos pesos iguais que somam 1.

\subsection*{Restrições práticas que mantêm a estrutura do questionário}
\begin{itemize}[leftmargin=1.5em]
  \item \textbf{Apenas colunas ajustáveis:} uma coluna (feature) $j$ só é ajustada se houver ao menos um valor $X_{:,j}>0$. Colunas com $X_{:,j}=0$ para todas as linhas são \emph{congeladas} e mantidas exatamente como no $W_0$ original (não são zeradas nem alteradas).
  \item \textbf{Limites dos pesos:} para colunas ajustáveis, projetamos $W$ no intervalo $[\,\varepsilon,\,1\,]$ (com $\varepsilon$ pequeno).
  \item \textbf{Proximidade ao original:} penalizamos desvios em relação a $W_0$ (pesos originais) via \emph{Elastic Net} (combinação de L1 e L2).
\end{itemize}

\subsection*{Função de perda e otimização}
Minimizamos a entropia cruzada média entre $P$ e $Y$ (somente nas classes clínicas tradicionais) com penalização elástica em torno de $W_0$:
\begin{equation}
  \mathcal{L}(W) \;=\; -\frac{1}{n}\sum_{i=1}^{n}\sum_{k=1}^{K} Y_{ik}\,\log P_{ik}
  \;+\; \lambda_1\lVert W-W_0 \rVert_1 \;+\; \lambda_2\lVert W-W_0 \rVert_2^2 \,.
\end{equation}
Após cada passo de gradiente/proximal, aplicamos a projeção
\begin{equation}
  \Pi(W) \;=\; 
  \begin{cases}
    W_{j,:} \leftarrow W_{0\,j,:}, & \text{se } \max_i X_{ij}=0 \quad (\text{coluna congelada})\\[2pt]
    W_{j,:} \leftarrow \mathrm{clip}(W_{j,:}, \varepsilon, 1), & \text{caso contrário (ajustável).}
  \end{cases}
\end{equation}
Paramos quando não há melhora do macro top-3 (ou ao atingir uma meta).

\section{Classe \textit{Sem Transtorno} por regra de ``baixa ativação''}
As features medem \emph{presença} de sintomas; por isso, ao invés de criar uma 12ª coluna em $W$ para o ``normal'', introduzimos a classe \textbf{Sem Transtorno} por uma \textbf{regra} aplicada sobre as probabilidades já calculadas $P$ (com $K=11$ classes).

\subsection*{Quando a regra aciona?}
Para cada linha, definimos:
\begin{align}
  p^{(1)} &= \max_k P_k \quad (\text{probabilidade top-1}),\\
  p^{(2)} &= \text{segunda maior probabilidade},\\
  \Delta  &= p^{(1)} - p^{(2)} \quad (\text{margem}).
\end{align}
A regra aciona se
\begin{equation}
  p^{(1)} < T_1 \quad \text{e} \quad \Delta < T_2 \,,
\end{equation}
sinalizando que não há classe fortemente ativada \emph{e} há pouca separação em relação à segunda colocada.

\subsection*{Como atribuímos probabilidade ao \textit{Sem Transtorno}?}
Se a regra aciona, alocamos uma fração $\gamma$ à nova classe e reescalamos as demais probabilidades para manter soma $1$.
Dado $P\in\mathbb{R}^K$,
\begin{equation}
  P' \;=\;
  \begin{cases}
    \big[(1-\gamma)P;\ \gamma\big], & \text{se aciona}\\[2pt]
    \big[P;\ 0\big], & \text{se não aciona}
  \end{cases}
\end{equation}
onde $P' \in \mathbb{R}^{K+1}$ inclui a posição extra para \emph{Sem Transtorno}.
O ranking (top-1/top-2/top-3) passa a considerar $K{+}1$ classes.

\subsection*{Como escolhemos $T_1$, $T_2$ e $\gamma$? (busca em grade)}
Avaliamos diversas combinações $(T_1,T_2,\gamma)$ e escolhemos a que maximiza o \textbf{macro top-3} \emph{incluindo} a classe \emph{Sem Transtorno}.
Para isso, mapeamos respostas ``não'' no \emph{Alvo} para este novo rótulo e computamos a taxa de acerto por classe (aparece no top-3) e depois a média entre classes com suporte $>0$.

\paragraph{Exemplo numérico.}
Suponha $P=[0{,}24,\ 0{,}22,\ 0{,}18,\ \ldots]$, então $p^{(1)}=0{,}24$, $p^{(2)}=0{,}22$, $\Delta=0{,}02$.
Se a melhor tripla for $T_1{=}0{,}30$, $T_2{=}0{,}06$, $\gamma{=}0{,}50$, a regra aciona e obtemos
\begin{equation}
  P'=\big[\,0{,}12,\ 0{,}11,\ 0{,}09,\ \ldots,\ 0{,}50\,\big],
\end{equation}
onde $0{,}50$ é $p_{\mathrm{Sem\ Transtorno}}$ e as demais probabilidades foram multiplicadas por $(1{-}\gamma)=0{,}5$.

\section{Leitura das planilhas geradas}
\begin{itemize}[leftmargin=1.5em]
  \item \textbf{Pontuação\_Tunada:} matriz de pesos $W$ após o ajuste (colunas congeladas iguais ao $W_0$, ajustáveis em $[\varepsilon,1]$).
  \item \textbf{Resultado\_Heuristica\_Tunada:} probabilidades $p\_{\langle classe\rangle}$, a coluna $p\_{\mathrm{Sem\ Transtorno}}$ e o ranking top-1/top-2/top-3 com as $K{+}1$ classes.
  \item \textbf{Metricas\_Heuristica\_Tunada:} macro top-3 (incluindo \emph{Sem Transtorno}) e tabela por classe (taxa de top-3, suporte).
  \item \textbf{Regras\_Normal:} valores aprendidos $(T_1,T_2,\gamma)$, taxa de acionamento e métrica final.
\end{itemize}

\paragraph{Observação.} Por configuração, podemos também exigir que uma coluna só seja ajustável se, além de $X_{:,j}$ ter algum valor $>0$, a coluna correspondente em $W_0$ possuir ao menos um peso diferente de zero (mantendo a \emph{sparsidade} original).

\end{document}